\chapter{缓存侧信道攻击防御方法研究}

\section{现有防御方法综述}\label{sec:defense_survey}

针对缓存侧信道攻击的防御研究已经积累了较多成果~\cite{ge2018survey,lou2022survey}。现有方法大体可分为软件层面和硬件/系统层面两类。本节简要梳理各类方法的原理与局限，并分析其在 AMD SEV-SNP 环境中对跨域缓存一致性逐出信道与物理内存竞争信道的适用性。

软件层面的防御主要从密码实现入手，尽量减少侧信道信号的来源。常数时间实现将密码算法改写为访问模式与数据值无关的形式。对于 AES，这通常意味着不再采用 T-Table 查表路径~\cite{Bernstein2005CachetimingAO,osvik2006cache}。更常见的做法是启用 AES-NI 等硬件密码加速指令，将 AES 运算移入专用执行单元，避免产生依赖密钥的内存访问~\cite{tromer2010efficient}。对 RSA 等非对称算法，可采用 Montgomery ladder 等常数时间指数运算替代含数据依赖分支的 Square-and-Multiply 实现~\cite{kocher1996timing}。此外，随机延迟或随机额外访问可以扰动单次观测，但难以消除统计意义上的信息泄露，攻击者仍可能通过重复观测恢复信号~\cite{lou2022survey}。这类软件防御通常需要来宾侧软件配合，宿主机难以单方面强制执行。

硬件与系统层面的防御主要是在共享资源上建立隔离边界。Intel CAT 和 AMD PQoS 等缓存资源管理技术可将末级缓存（LLC）划分为互不重叠的路集合，使机密虚拟机占用特定 LLC 分区，从而抑制依赖 LLC 共享的 Prime+Probe、Flush+Reload 等攻击~\cite{intel_cat_spec,amd_pqos_ppr,liu2016catalyst}。缓存替换策略随机化（如 CEASER~\cite{qureshi2018ceaser}）通过地址置换降低冲突集的可预测性，但需要处理器微码或设计层面的支持。基于性能监控单元（PMU）的异常检测则通过监控 LLC 缺失率等指标识别高频探测行为~\cite{gast2025counterseveillance}。这类方法属于被动防御，对低频慢速攻击的灵敏度有限，检测阈值也不容易在误报和漏报之间取得平衡。

LLC 分区对本文所研究的两类信道作用并不相同。隔离 LLC 路集合后，宿主机仍可以宿主侧地址标签访问目标物理页，C-bit 不匹配导致的跨域一致性逐出仍可能发生，因此通常只能起到部分缓解作用。另一方面，NOCACHE 探测绕过处理器缓存直接与内存控制器交互，LLC 分区对物理内存竞争信道的抑制较为有限。各方案的综合对比见表~\ref{tab:defense_comparison}。

来宾侧启用 AES-NI 对两类信道都有较直接的抑制作用，但这一路径依赖来宾软件配合，宿主机无法强制要求。在需要由宿主机单方面实施防御的场景下，LLC 分区仍是较可行的选项。不过，它对内存竞争信道的覆盖不足，而这一点也对应了本文攻击系统中竞争特征的提取路径。基于这一情况，下一节在 AMD PQoS 的基础上讨论一种面向机密虚拟机的增强型缓存隔离防御方案，并分析其对两类信道的缓解能力。

\begin{table}[htbp]
  \centering
  \caption{现有防御方法在 SEV-SNP 环境中的对比}\label{tab:defense_comparison}
  \begin{tabular}{lcccc}
    \toprule
    防御方法 & 部署方 & 对一致性逐出信道 & 对内存竞争信道 & SEV-SNP 可部署性 \\
    \midrule
    常数时间实现（AES-NI） & 来宾 & 可消除 & 可消除 & 需来宾配合 \\
    常数时间实现（软件）   & 来宾 & 可缓解 & 可缓解 & 需来宾配合 \\
    随机化噪声注入         & 来宾 & 增加攻击难度 & 增加攻击难度 & 需来宾配合 \\
    LLC 分区（CAT/PQoS）  & 宿主机 & 部分缓解 & 作用有限 & 需可信管理程序 \\
    缓存随机化             & 硬件  & 作用有限 & 作用有限 & 需硬件支持 \\
    PMU 异常检测           & 宿主机 & 可检测部分场景 & 可检测部分场景 & 可在宿主侧实施 \\
    \bottomrule
  \end{tabular}
\end{table}

\section{基于 AMD PQoS 的增强型缓存隔离防御方案}
\label{sec:pqos_defense}

本节讨论一种面向 SEV-SNP 机密虚拟机的两层缓存隔离方案。方案以 AMD PQoS 提供的缓存资源控制能力为基础，对两类信号分别采取不同的约束方式。LLC 路集合隔离主要作用于跨域缓存一致性逐出信道，MBA 内存带宽限速主要作用于物理内存竞争信道。与仅对 LLC 做路分区的传统方案相比，这一思路更接近两类信号各自的形成位置。

AMD EPYC 处理器的 Platform QoS 功能通过模型相关寄存器对末级缓存的路集合和内存带宽分配进行软件控制~\cite{amd_pqos_ppr}。每个逻辑处理器关联一个类服务标识符 CLOSID，系统为不同 CLOSID 配置 LLC 位掩码和带宽约束，从而控制该逻辑处理器可使用的缓存路集合及其内存访问份额。Linux 内核通过 \texttt{resctrl} 文件系统暴露这些能力，为方案实施提供了现成接口。基于这一机制，本文将方案分为两个层次。

第一层是 LLC 路集合隔离。实验平台 AMD EPYC 7763 的 LLC 共 32 路。本文将 LLC 路分配掩码等分为两组，低半部分路集合（way mask \texttt{0x00FF}）分配给宿主机线程所属的 CLOSID，高半部分（\texttt{0xFF00}）分配给机密虚拟机 vCPU 所属的 CLOSID。启用后，宿主机线程与来宾 vCPU 在 LLC 中占用互不重叠的路集合，无法在同一路上相互驱逐。对于跨域缓存一致性逐出信道，这一层的作用在于减少宿主侧 \texttt{CACHEABLE} 探测与来宾访问之间形成缓存共享的机会。隔离路集合后，宿主侧副本与来宾侧副本落入不同的缓存路，相关冲突条件被削弱。对其他依赖 LLC 共享的攻击路径，如 Prime+Probe 和 Flush+Reload，该层通常也有缓解作用。对于 NOCACHE 探测，这一层的影响较小，因为后者绕过 LLC，直接与内存控制器交互。

第二层是 MBA 内存带宽限速。物理内存竞争信号的形成依赖宿主机在短测量窗口内发起高频 \texttt{NOCACHE} 突发访问，并与来宾的 DRAM 请求争用内存控制器带宽。对宿主机线程所属的 CLOSID 施加带宽上限后，宿主机在单位时间内可发起的有效内存访问减少，突发探测的带宽密度随之下降，竞争窗口内的延迟差异一般会减弱，攻击者需要更多轮观测才能积累可用的统计信号。该层作用于 DRAM 访问层，可作为 LLC 分区之外的补充约束。

两层机制需要由可信配置路径统一控制。在 SEV-SNP 的威胁模型下，若将 PQoS 路分配和 MBA 参数设置暴露给 QEMU 用户态，攻击者仍可能在防御启用后重新配置资源控制寄存器，从而削弱防御效果。本文将配置权收敛到宿主机内核中的 KVM 扩展模块。该模块在机密虚拟机启动阶段根据虚拟机标识自动完成 CLOSID 分配和 \texttt{resctrl} 控制组写入，配置完成后锁定相关路径，不向宿主机用户态暴露可重入的配置接口。资源隔离的建立与维持由 Hypervisor 内核层负责，QEMU 不直接控制防御参数。图~\ref{fig:pqos_defense}~给出了该方案的整体架构。

\begin{figure}[htbp]
  \centering
  \includegraphics[width=0.95\textwidth]{pic/pqos_defense.pdf}
  \caption{基于 AMD PQoS 的两层缓存隔离防御方案架构}
  \label{fig:pqos_defense}
\end{figure}

\section{防御原理实验验证}
\label{sec:defense_experiment}

受时间条件限制，本文未完成该防御方案的完整实现。本节通过两组原理性验证实验，分别评估 LLC 路集合隔离和 MBA 带宽限速对跨域缓存一致性逐出信号与物理内存竞争信号的缓解效果。实验环境与第~\ref{chap:channel}~章保持一致，均运行于 AMD EPYC 7763 平台，来宾为单 vCPU 的 SEV-SNP 机密虚拟机。

实验的目的是观察防御措施下两类信号是否仍然存在，以及其强度如何变化。测量期间，来宾虚拟机程序以固定节奏交替触发两类 AES 明文输入。A 组输入使 \texttt{Te0[0..7]} 被访问，B 组输入使 \texttt{Te0[128..255]} 被访问。两组访问对应 \texttt{Te0} 中不同的物理页内位置，其缓存行级访问模式并不相同。宿主线程在每次 AES 调用后立即探测目标缓存行并记录 \texttt{tsc\_delta}，同时记录当前样本所属的 A/B 组标签。若信号仍然存在，A/B 两组延迟分布通常会表现出一定差异；若防御起作用，这种差异应当缩小。一致性逐出实验使用 \texttt{CACHEABLE} 探测模式，burst 为 1 次。竞争信号实验使用 \texttt{NOCACHE} 探测模式，burst 为 16 次。每组各采集 1500 个样本。

\subsection{LLC 分区对一致性逐出信号的缓解效果}

启用 \texttt{resctrl} 进行 LLC 路集合隔离后，宿主组 LLC 路掩码配置为 \texttt{0x00FF}，CVM 组配置为 \texttt{0xFF00}，两组路集合互不重叠。实验结果（$n=1500$）为 A 组中位延迟 711 cycles，B 组中位延迟 711 cycles，两组中位数相同，信号分离值（$\text{separation\_median}$）为 0，信噪比（SNR）为 0。

与第~\ref{chap:channel}~章的基线测量相比，在无防御条件下，相同的 A/B 组访问模式在 \texttt{CACHEABLE} 探测下会产生延迟差异，中位延迟差超过 400 cycles，ROC-AUC = 0.999648。启用 LLC 路集合隔离后，两组分布在中位数层面不再表现出差异，SNR 降为 0。就当前 \texttt{CACHEABLE} 探测路径和实验设置而言，这表明一致性信号有所减弱。

\subsection{MBA 限速对物理内存竞争信号的缓解效果}

在 LLC 路集合隔离的基础上，对宿主组额外施加 20\% 的 MBA 内存带宽上限。实验结果（$n=1500$）为 A 组中位延迟 548.5 cycles，B 组中位延迟 549.0 cycles，中位数差为 0.5 cycles，SNR 降至 0.0084。

在无 MBA 限速的基线条件下，\texttt{NOCACHE} burst 探测可在竞争窗口内对两类访问模式产生一定的延迟差异，相关结果见第~\ref{chap:channel}~章中的竞争信号测量。施加 20\% 带宽上限后，突发探测的有效窗口内两组延迟接近重合，信噪比接近 0。就当前测量框架而言，这说明物理内存竞争信号已经下降到较难直接利用的水平。

两组实验结果表明，LLC 路集合隔离与 MBA 带宽限速分别对两类信号路径具有缓解作用。两层机制均通过 \texttt{resctrl} 在宿主机侧实施，无需修改来宾软件，与 SEV-SNP 的威胁模型基本一致。

\subsection{防御机制的性能损耗}

除信号抑制效果外，防御方案还需要满足较低的运行时开销。为此，本文进一步在来宾内部运行纯 AES benchmark，对防御启用前后的执行时间进行比较,记录总执行时间、总周期数以及每次调用的平均开销。实验结果表明，两类宿主侧防御对来宾内部 AES 运算的影响都较小。在无防御的 baseline 条件下，guest 内部 \texttt{AES\_encrypt} 的中位执行开销为 89.275699 ns/call，对应 218.725039 cycles/call。启用 LLC 路集合隔离后，中位执行开销为 89.232517 ns/call，对应的相对变化为 $-0.048\%$。该差值处于重复实验的波动范围内，可认为未引入可观测的额外开销。启用 20\% MBA 带宽限制后，中位执行开销为 89.278697 ns/call，相对 baseline 的变化为 $+0.0034\%$，同样接近于零。

从性能角度看，当前 \texttt{resctrl} 配置主要约束的是宿主线程可使用的 LLC 路集合与内存带宽，而不会显著改变来宾内部 AES 计算路径本身。因此，在本文所采用的工作负载和平台配置下，LLC 分区与 MBA 带宽限制均表现出较低的性能代价。结合前述信号抑制结果，可以认为这两类宿主侧资源控制机制在当前实验平台上具有较好的防御收益与性能平衡。
